\subsection*{一步一步建立直观}

\paragraph{1. 在输出中追踪一根小输入箭头}

把观察点固定，在输入轴上画一段小位移 $h$，再比较输出变化。曲线通常不是线性变换，但在足够小的尺度上，输出变化可由一个固定倍率近似。这个倍率就是导数。右图的割线把有限位移联系起来，切线显示局部替代模型。



\paragraph{2. 缩小尺度，分离主项与误差}

局部线性化应写成 $f(a+h)-f(a)=f'(a)h+o(h)$。余项不仅要趋零，还要相对于输入尺度 $h$ 更小。平方函数的剩余项是 $h^2$，因此除以 $h$ 后消失。这个相对误差条件才支持把导数视作真正的一阶变换。



\paragraph{3. 从一个倍率走向矩阵}

移动观察点，局部倍率随位置更新。在多变量情形，小位移有多个方向，导数便成为作用在位移向量上的线性映射，用雅可比矩阵表示。本图只展示一维版本；链式法则在多维中相应变成矩阵复合，并须保持输入输出空间的次序。


